% !TEX root = ../main.tex

\chapter{Introduction}

Domain theory is a powerful tool for dealing with partial functions in the denotational semantics of programming languages. The classical approach assigns each type a domain, which is a poset with directed suprema, and sees functions between these types as homomorphisms between domains \cite{scott_domains_1982,scott_outline_1970}. Meanwhile, the classical approach does not immediately scale to different computation models with different strengths of computability, or models like concurrent computations \cite{nygaard_domain_2004,winskel_powerdomains_1985}.

Synthetic domain theory (SDT) is a synthetic and axiomatic approach to domain theory \cite{fiore_axiomatic_1996,hyland_first_1991,reus_general_1997}. In SDT, the domains and orders are not introduced a priori. Instead, it constructs the theory in a topos \cite{mac_lane_sheaves_1994} with a classifier that classifies a certain class of subsets \cite{hyland_first_1991}. The illustration of partial functions and the construction of domains are then derived from this classifier.

However, defining the orders inside SDT is a non-trivial topic and leads to different choices of axioms that yield different models of SDT \cite{van_oosten_axioms_2000}. One promising approach is inspired by the homotopy type theory (HoTT) \cite{the_univalent_foundations_program_homotopy_2013}. HoTT gives a type theory that treats identity types as paths and uses algebraic topological constructions to give important properties related to identity types, including congruence, functional extensionality, and univalence. Due to the nature of identity types and homotopy, HoTT can be seen as a theory of $\infty$-groupoids.

Inspired by this, Riehl and Schulman gave a variant HoTT known as directed homotopy type theory (directed HoTT) and constructed a synthetic theory of $\infty$-categories \cite{riehl_type_2023}. To restrict a type in directed HoTT to behave like a $\infty$-category, the type must satisfy two completeness properties called Segal and Rezk. By seeing a poset as a thin $1$-category and taking the interval type as the subobject classifier, directed HoTT can therefore be used to axiomatise SDT, and hence we need to find proper axioms to imply that the interval itself is a domain, which includes showing that the interval is Segal and Rezk.

Phoa principle is one property about the classifier \cite{phoa_domain_1991} that is proven to be enough to imply the Segal and Rezk properties for the interval type along with other properties that altogether show that the interval is a poset \cite{sterling_domains_2025}. The Phoa principle has also been generalised to higher-dimensional simplices, and such generalisation is shown to be equivalent to the original Phoa principle \cite{pugh_when_2025}.

This thesis proposes a further generalisation of the Phoa principle and provides verified proofs in Cubical Agda \cite{vezzosi_cubical_2019}. Specifically, the thesis contains the following works:
\begin{enumerate}
    \item Rephrasing the higher Phoa principle by introducing a dual version of simplices.
    \item Generalising the higher Phoa principle to the transfinite cases based on the rephrasing.
    \item Introducing the concept of sobriomorphisms and revealing a possible connection between the Phoa principle and the chain-completeness, a property that makes a synthetic poset a domain.
    \item Axiomatising the interval type in Cubical Agda and proving the classical, higher, and transfinite Phoa principle with concrete constructions in Cubical Agda.
\end{enumerate}

The thesis starts with the algebraic structures on the interval type and illustrates the connection between the partial map classifier and the simplices in Chapter \ref{chapter-interval}. In Chapter \ref{chapter-phoa}, we give a review on the classical and higher Phoa principle with a rephrasing based on the dual simplices and sobriomorphisms. Then in Chapter \ref{chapter-omega}, we generalise the Phoa principle to the transfinite case that relates the simplices and the initial and final (co)algebras defined by the partial map classifier. Finally, in Chapter \ref{chapter-topology}, we make a brief investigation on the relation between the Phoa principle and the completeness properties involved in synthetic domains, and propose a conjectural completeness property that may possibly unify the Segal and the chain completeness. Several theorems and lemmata in the text are accompanied by an Agda object name written in monospace fonts. These refer to the corresponding formalisation code in the Agda codebase.
